Import 10 mln biletów najszybciej wykonuje Valkey
(\lval{valkey}{10000000}{load}\,s), potem PostgreSQL
(\lval{postgres}{10000000}{load}\,s, binarny \texttt{COPY}), MariaDB
(\lval{mariadb}{10000000}{load}\,s, wielowierszowe \texttt{INSERT}) i~MongoDB
(\lval{mongo}{10000000}{load}\,s). MongoDB płaci za \emph{write concern}
\texttt{majority} z~dziennikiem i~za budowę dokumentów BSON z~osadzonymi
tablicami. Czas importu rośnie z~rozmiarem w~przybliżeniu liniowo we
wszystkich bazach.

Tworzenie 15 indeksów na 10 mln biletów trwa w~PostgreSQL
\lval{postgres}{10000000}{indexes_on}\,s, w~MariaDB
\lval{mariadb}{10000000}{indexes_on}\,s, a~w~MongoDB
\lval{mongo}{10000000}{indexes_on}\,s. W~Valkey ,,indeksy'' buduje aplikacja:
musi przeczytać przez \texttt{SCAN} wszystkie 1,67 mln rezerwacji i~zapisać
zbiory posortowane, co zajmuje \lval{valkey}{10000000}{indexes_on}\,s.

Najmniej miejsca zajmuje MongoDB dzięki kompresji WiredTiger (Snappy):
\sval{mongo}{10000000}{off}\,MB bez indeksów wobec
\sval{postgres}{10000000}{off}\,MB w~PostgreSQL i~\sval{mariadb}{10000000}{off}\,MB
w~MariaDB. Valkey przy 10 mln biletów zajmuje \sval{valkey}{10000000}{off}\,MB
pamięci RAM. Cały zbiór musi się w~niej zmieścić, co jest praktyczną granicą
tego modelu. Wartości ,,z~indeksami'' obejmują też wiersze dopisane przez
scenariusze z~pierwszego przebiegu (np. C3 importuje \num{10000} biletów
w~każdym powtórzeniu).
